% Options for packages loaded elsewhere
\PassOptionsToPackage{unicode}{hyperref}
\PassOptionsToPackage{hyphens}{url}
%
\documentclass[
]{article}
\usepackage{amsmath,amssymb}
\usepackage{lmodern}
\usepackage{iftex}
\ifPDFTeX
  \usepackage[T1]{fontenc}
  \usepackage[utf8]{inputenc}
  \usepackage{textcomp} % provide euro and other symbols
\else % if luatex or xetex
  \usepackage{unicode-math}
  \defaultfontfeatures{Scale=MatchLowercase}
  \defaultfontfeatures[\rmfamily]{Ligatures=TeX,Scale=1}
\fi
% Use upquote if available, for straight quotes in verbatim environments
\IfFileExists{upquote.sty}{\usepackage{upquote}}{}
\IfFileExists{microtype.sty}{% use microtype if available
  \usepackage[]{microtype}
  \UseMicrotypeSet[protrusion]{basicmath} % disable protrusion for tt fonts
}{}
\makeatletter
\@ifundefined{KOMAClassName}{% if non-KOMA class
  \IfFileExists{parskip.sty}{%
    \usepackage{parskip}
  }{% else
    \setlength{\parindent}{0pt}
    \setlength{\parskip}{6pt plus 2pt minus 1pt}}
}{% if KOMA class
  \KOMAoptions{parskip=half}}
\makeatother
\usepackage{xcolor}
\setlength{\emergencystretch}{3em} % prevent overfull lines
\providecommand{\tightlist}{%
  \setlength{\itemsep}{0pt}\setlength{\parskip}{0pt}}
\setcounter{secnumdepth}{-\maxdimen} % remove section numbering
\ifLuaTeX
  \usepackage{selnolig}  % disable illegal ligatures
\fi
\IfFileExists{bookmark.sty}{\usepackage{bookmark}}{\usepackage{hyperref}}
\IfFileExists{xurl.sty}{\usepackage{xurl}}{} % add URL line breaks if available
\urlstyle{same} % disable monospaced font for URLs
\hypersetup{
  hidelinks,
  pdfcreator={LaTeX via pandoc}}

\author{Dietmar Gerald Schrausser}
\date{09.05.2026}


\begin{document}

\hypertarget{foliumblat-iir}{%
\section{Folium/Blat IIr}\label{foliumblat-iir}}

Eine zeitgenössische deutsche Übersetzung wird präsentiert, die auf
einem Vergleich der bestehenden lateinischen, frühneuhochdeutschen
(ENHG) und englischen (Übersetzungs-) Texte basiert, um ein besseres
Verständnis zu ermöglichen. Insbesondere da der ENHG-Text im Vergleich
zum modernen Hochdeutschen schwer verständlich ist, da dieser fundierte
Deutschkenntnisse (als Muttersprache) sowie Kenntnisse verschiedener
Dialekte voraussetzt.

\hypertarget{ohne-uxfcberschrift.}{%
\subsection{Ohne Überschrift.}\label{ohne-uxfcberschrift.}}

\begin{quote}
IN principio creauit deus celum et terra{[}m{]}. Terra autem erat inanis
{[}et{]} vacua. {[}et{]} tenebrae era{[}n{]}t sup{[}er{]} faciem abissi.
et sp{[}irit{]}us d{[}omi{]}ni ferebat{[}ur{]} sup{[}er{]} aq{[}ua{]}s.
\end{quote}

\textbf{Moses}, der göttliche Prophet und Historiker, welcher nahezu 700
Jahre vor dem Trojanischen Krieg gelebt hat, lehrt, dass Gott der Planer
und Schöpfer der Dinge, um dieses Werk zu vollbringen zuallererst den
Himmel schuf und ihn überordnete, als Sitz des Schöpfergottes selbst.
Dann gründete er die Erde und ordnete sie dem Himmel unter. Er hüllte
die Erde in Finsternis, d.h. diese ist ohne eigene
Erkenntnis\footnote{`dan{[}n{]} sie begreüft durch sich selbs nichtzit'.},
es sei denn sie empfängt das Licht vom Himmel. Dort setzte er das ewige
Licht, die himmlischen Geister und das ewige Leben ein, der Erde
hingegen gab er Finsternis, die irdischen Geister und den
Tod.\footnote{von Lactantius
  (\href{https://books.google.com/books?id=qGsnghsi3uUC}{1497}, Buch 2,
  Kapitel 9/10).}

Dadurch, dass Moses eine göttliche \emph{Schöpfung} beschreibt, zeigt er
die drei Irrtümer\footnote{Bezeichnet die klassischen philosophischen
  Ansichten über die Schöpfung und die Ewigkeit des Universums, die der
  Genesis widersprachen.} des \textbf{Platon}, \textbf{Aristoteles} und
\textbf{Epikurs} auf. Denn Platon war der Ansicht, dass (1) Gott, (2)
die \emph{Ideen}\footnote{`ydeas' und `die vorpildnus oder gestaltnus
  seiner geschöpff', Ideen (Ydeas) sind die ewigen, vollkommenen
  Archetypen im platonischen Denken.} und Yle von Ewigkeit gewesen sind,
die Welt (3) grundsätzlich aus (diesem) Yle gemacht worden sei.
\emph{Yle} bezeichnen die Griechen die erste formlose Masse aus der
\emph{alles geschaffen} wurde; die \emph{sichtbaren} Dinge, die aus
harmonisierenden Elementen geformt sind. Oder auch, aus Materie und
Form, dem feinsten Staub\footnote{`de athomis' und `de{[}m{]} aller
  dynnisten staub'.} der im Sonnenlicht glitzert, gemacht.\footnote{wird
  der `Summa super Sententias' (Petrus,
  \href{https://doi.org/10.3931/e-rara-18241}{1484}, Buch II, Abschnitt
  1) zugeschrieben, behandelt Platons Ansicht über die ewige, ungeformte
  Materie, Aristoteles' Behauptung einer ewigen Welt und Epikurs Theorie
  der atomaren Kollision.}

(ad 3) Dennoch\footnote{`t{[}ame{]}n'.} erschuf Gott die Welt ohne
verfügbares und zuvor vorbereitetes Material; denn \emph{er} ist als der
klügste\footnote{`prude{[}n{]}tissim{[}us{]}'.} und weiseste Schöpfer
anzusehen, bevor er das Werk der Welt vollführte. In ihm selbst war der
Quell vollkommener und vollbrachter Güte, damit\footnote{`vt'.} daraus
Gutes gleich einem \emph{Strom}\footnote{`tam{[}que{]}m riu{[}us{]}' und
  `als ein pach'.} hervorgebracht wird.\footnote{Dient als
  christlich-neuplatonisches Argument für die Schöpfung aus dem Nichts
  (\emph{creatio ex \textbf{nihilo}}) und stellt sie griechischen
  philosophischen Ansichten gegenüber. Es betont, dass \emph{Gott} nicht
  zufällig, sondern \emph{bewusst} (`excogitandum') handelte (vgl.
  Trismegistu,
  \href{https://gdz.sub.uni-goettingen.de/id/PPN858183994}{1471}).}

(ad 2) Von allen Wesen erschuf er ursprünglich die
\emph{Engel}\footnote{Die intellektuelle Welt, s. fol.~1r `quem theologi
  angelicum: philosophi: autem intellectualem vocant'.} aus dem was
\emph{Nicht} (\textbf{\emph{n{[}{]}}}) ist\footnote{`ex eo q{[}uo{]}d
  n{[}{]} e{[}st{]}' und `auß de{[}n{]} das nicht ist'.}. So ist
\emph{er} durch die \emph{Ewigkeit} mächtig und durch die Stärke dieser
immensen Macht (welche grenzenlos und unermesslich ist\footnote{`q{[}uam{]}
  fine ac mo{[}do{]} caret' und `die des ends vn{[}d{]} der maß
  mangelt'.}) \emph{gleich} dem Leben des
Schöpfers\footnote{`sic{[}ut{]} vita facturis' vs.~`als das lebe{[}n{]}
  des Schöpffers'.}.\footnote{erörtert das christliche theologische
  Konzept der Schöpfung aus dem Nichts (\emph{creatio ex nihilo}).}
Daher mag es verwunderlich sein, dass zur Schöpfung der Welt im
Vorhinein eine Roh-Materie\footnote{`materia{[}m{]} de q{[}ua{]}
  faceret' und `ein materi daraus er machet'.} aus dem \emph{Nicht}
vorzubereitet ist.\footnote{vgl. Lactantius
  (\href{https://books.google.com/books?id=qGsnghsi3uUC}{1497}, Buch II
  8.8), er argumentierte gegen antike griechische Philosophien, die
  behaupteten, Gott habe die Welt aus \emph{bereits existierender}
  Materie erschaffen.} So haben es wohl auch die Sarazenen \footnote{aus
  den `Etymologiae' (De Seville,
  \href{https://books.google.com/books?id=wb8Z_vlSyJwC}{1802}, Buch
  VIII, Kapitel xi, \emph{De diis gentium}), geschrieben im
  \emph{frühen} 7. Jahrhundert, um 610 n.~Chr. Es kann in diesem Fall
  davon ausgegangen werden, dass auf die \emph{Sarazenen} aus dem 4.
  Jahrhundert referiert wird, nomadische Stämme der syrisch-arabischen
  Wüste und es sich daher nicht um einen Bezug zum \emph{Islam} handelt.
  Vielmehr scheint es sich auf die \emph{weichenstellenden} Ereignisse
  um 378 n.~Chr. zu beziehen, im Zusammenhang mit der verheerenden
  Niederlage des Weströmischen Reiches gegen die Goten bei Adrianopel
  und der Beteiligung der \emph{Sarazenen} unter \emph{Mavia}
  (\emph{Mawiyya}), einer arabischen Herrscherin, die kurz zuvor im
  südlichen Syrien einen Aufstand gegen das Römische Reich angeführt und
  anschließend \emph{die Seiten gewechselt hatte}. Interessanterweise
  fällt auch die Geschichte vom \emph{Heiligen} \textbf{\emph{Moses}}
  \emph{dem Bischof der Araber} in diesen Kontext - nicht zu verwechseln
  mit dem gleichnamigen \textbf{\emph{Moses}} \emph{dem Schwarzen},
  einem bekehrten Verbrecher, der zur selben Zeit in der selben Region
  agierte (ebenfalls bekannt unter den Namen \emph{Moses der Starke},
  \emph{Moses der Abessinier}, \emph{Moses der Inder}, weiters
  \emph{Abba Moses der Räuber}, \emph{Moses Murin} oder \emph{Moses der
  Äthiopier} etc.).} verstanden als sie sagten, dass die Engel von Gott
aus der \emph{Finsternis} ins Licht geführt und mit ewiger Freude
erfüllt worden seien. Jedoch vergaßen einige von ihnen ihren göttliche
Charkter\footnote{`indolis' und `einpildung'.}, wandten sich durch
\emph{eigene} Abkehr\footnote{`verkerung', pervertierung.} vom Guten zu
Übel\footnote{`ex bono p{[}er{]} se malu{[}m{]} effectu{[}m{]}' und
  `vo{[}m{]} gůtten zum ubel getretten'.} und wurden zu Teufeln (was wir
Verbrecher\footnote{`greci diabolu{[}m{]} appelant: nos
  crimi{[}n{]}atore{[}m{]} vocam{[}us{]}'.} nennen).\footnote{aus dem
  `Heptaplus' (Mirandola \& Mirandola,
  \href{https://doi.org/10.3931/e-rara-61217}{1601}, Buch 1, Kapitel 1).}

(ad 1) Die Erde war leer, das heißt (wie \textbf{Hieronymus} und
\textbf{die Siebzig}\footnote{`septinge{[}n{]}ta', 700 vs.~`die .lxx.
  auslege{[}r{]}'.} es deuten) unsichtbar und nicht
zusammengefügt\footnote{`i{[}n{]}co{[}m{]}posita' und `vnzesamen
  gefügt', bezieht sich auf die lateinischen Übersetzungen des
  hebräischen `tohu wa-bohu' (formlos und leer) aus Genesis 1,2 (vgl.
  Jiménez de Cisneros,
  \href{https://doi.org/10.3931/e-rara-46695}{1517}).} und aufgrund
dieses konfusen\footnote{`{[}con{]}fusio{[}n{]}e' und `zestrewlichkeit'.}
Zustandes \emph{Abgrund} (\emph{Abyssum}) genannt, was die Griechen
\emph{Chaos} nennen. Das ist\footnote{`i{[}d est{]}'.}
\emph{dreidimensionale} Materie\footnote{`materia{[}m{]} trino
  dime{[}n{]}su' und `materi mit driueltiger ermessung'.}, in grösste
Tiefe ausgedehnt\footnote{`i{[}n{]} altissimas profunditates extensam'.}.

Auch Ovid erinnert sich daran:

\begin{quote}
Ante mare {[}et{]} terras {[}et{]} q{[}uo{]}d tegit o{[}mn{]}ia
celu{[}m{]}. Un{[}us{]} erat toto nature vult{[}us{]} in orbe.
Que{[}m{]} dixere chaos rudis indigestaq{[}ue{]} moles. Nec
q{[}ui{]}cq{[}uam{]} nisi po{[}n{]}dus iners: {[}con{]}gestaq{[}ue{]}
eode{[}m{]}. No{[}n{]} bene iu{[}n{]}ctaru{[}m{]} discordia femina
reru{[}m{]}. Nullus ad huc mu{[}n{]}do p{[}re{]}bebat lumina titan.
\end{quote}

So schwebte der Geist des Herrn (als Werkzeug göttlicher Kunst) über den
Wassern, wie der \emph{Wille} eines Baumeisters, der alles organisiert
was zu tun ist. Dadurch sind die Werke Gottes vollkommen und drückt sich
die Schöpfung in der \emph{Sechserzahl}\footnote{`senario numero' und
  `in sechser zall'.} aus, deren Bestandteile 1, 2 und 3 sind, welche
sich zu einem Dreieck erheben\footnote{`que in trigonu{[}m{]}
  surga{[}n{]}t'.}. Moses widmet in seinem Bericht über die sechs Tage
den ersten Tag (i) der Schöpfung, den zweiten und dritten Tag (ii) der
Ordnung und Gestaltung und (iii) die übrigen Tage der Zierde der Welt.
\footnote{aus dem `Heptaplus' (Mirandola \& Mirandola,
  \href{https://doi.org/10.3931/e-rara-61217}{1601}), in der klassischen
  und mittelalterlichen Zahlentheorie ist 6 die erste \emph{vollkommene
  Zahl}, weil sie gleich der Summe ihrer echten Teiler ist (1+2+3=6).}

\hypertarget{references}{%
\subsection{References}\label{references}}

De Seville, I. (1802). \emph{S. Isidori Hispalensis Episcopi Hispaniarum
Doctoris Opera Omnia}. Edited by Arevalo, F. Romae: Apvd Antonivm
Fvlgonivm. \url{https://books.google.com/books?id=wb8Z_vlSyJwC}

Jiménez de Cisneros, F. (1517). \emph{Biblia Polyglotta Complutensis}.
Complutum: Arnaldo Guillén de Brocar.
\url{https://doi.org/10.3931/e-rara-46695}

Lactantius, L.C.F. (1497). \emph{Divinae Institutiones}. Venetiis: Simon
Bevilaqua. \url{https://books.google.com/books?id=qGsnghsi3uUC}

Mirandola, G., \& Mirandola, G. F. (1601). \emph{Ioannis Pici,
Mirandulae \ldots{} opera quae extant omnia \ldots{}} Basileae: per
Sebastianum Henricpetri. \url{https://doi.org/10.3931/e-rara-61217}

Petrus. (1484). \emph{Sententiarum Libri IV}. Basel.
\url{https://doi.org/10.3931/e-rara-18241}

Trismegistu, H. (1471). \emph{Mercurii Trismegisti Liber de Potestate Et
Sapientia Dei e Graeco in Latinum Traductus a Marsilio Ficino Pimander
Incipit}. Edited by Ficinus, M. Treviso.
\url{https://gdz.sub.uni-goettingen.de/id/PPN858183994}

\end{document}
